\emph{Setting.} Throughout, the projected coherent vector is read as the $\sqrt{m}$ vector (its circulant reading); every statement below is conditional on that reading. By Proposition \claimref{synthesis.coherent-face-vector-real-slices}(iii) the projected coherent face vector is $p_{k} = |\kappa|(1 + \lambda\cos(\delta + \tfrac{2\pi k}{3}))$ with $\delta = \chi$ on the Euclidean slice and $\delta = \chi \mp \tfrac{\pi}{2}$ on the Lorentzian one (by the orientation of the singlet coordinate), and by (iv) it satisfies the Koide relation exactly when $\lambda^{2} = 2$, which is the null condition of the equal-coefficient member of Proposition \claimref{families.koide-nullity-condition-one-member-invariant}. Drop the common factor and take $\lambda = \sqrt2$. Relabelling the families shifts $\delta$ by multiples of $\tfrac{2\pi}{3}$, reflecting the labels sends $\delta \mapsto -\delta$, and reversing the sign of the doublet part shifts $\delta$ by $\pi$; together these generate $\pm\delta + \tfrac{\pi}{3}\mathbb{Z}$, and with the slice choice the channel phase lies on $\pm\delta + \tfrac{\pi}{6}\mathbb{Z}$. Write $c = \cos\delta$, $s = \sin\delta$; then $\cos(\delta + \tfrac{2\pi}{3}) = -\tfrac{c}{2} - \tfrac{\sqrt3}{2}s$ and $\cos(\delta + \tfrac{4\pi}{3}) = -\tfrac{c}{2} + \tfrac{\sqrt3}{2}s$, so $r(\delta) = (p_{2}/p_{1})^{2}$ is evaluated without any value of $\pi$.

\emph{(i)} On $[0, \tfrac{\pi}{6}]$, $\delta + \tfrac{2\pi}{3}$ runs over $[\tfrac{2\pi}{3}, \tfrac{5\pi}{6}]$, where the cosine falls from $-\tfrac12$ to $-\tfrac{\sqrt3}{2}$ and crosses $-\tfrac{1}{\sqrt2}$ at $\delta = \tfrac{\pi}{12}$; so $p_{1} > 0$ exactly for $\delta < \tfrac{\pi}{12}$ (the positive root, as Proposition \claimref{synthesis.light-cone-koide-is-quarter-turn}(iv) selects it). On $[0, \tfrac{\pi}{12}]$, $\delta + \tfrac{4\pi}{3} \in [\tfrac{4\pi}{3}, \tfrac{17\pi}{12}]$, where the cosine rises from $-\tfrac12$, so $p_{2}$ increases and stays above $1 - \tfrac{\sqrt2}{2} > 0$, while $p_{1}$ decreases to $0$; hence $r = (p_{2}/p_{1})^{2}$ increases strictly from $1$ (at $\delta = 0$, $p_{1} = p_{2}$) to $\infty$, and $r(\delta) = r_{\mathrm{obs}}$ has exactly one solution there. The check evaluates $r$ in certified interval arithmetic: $\sin$ and $\cos$ of a rational in $(0, 1)$ by Taylor polynomials of degree $47$ with twice the first omitted term as error (the series alternates with decreasing terms), $\sqrt2$, $\sqrt3$ by integer square roots at $10^{-40}$, every interval end rounded outwards to $10^{-40}$. The window for $r_{\mathrm{obs}}$ is $((m_{\mu} - 3\sigma_{\mu})/(m_{e} + \epsilon_{e}), (m_{\mu} + 3\sigma_{\mu})/(m_{e} - \epsilon_{e}))$ at the stored values $m_{\mu} = 105.6583755$ MeV, $\sigma_{\mu} = 2.3 \times 10^{-6}$ MeV, $m_{e} = 0.51099895$ MeV, with $\epsilon_{e} = 10^{-9}$ MeV, which covers both the stored uncertainty and the store's eight-decimal cut of the CODATA value. The check certifies $r(0.22222204589) < r_{\min}$ and $r(0.22222204830) > r_{\max}$, that neither statement survives one step of $10^{-11}$ inwards (the window is tight at that resolution), the positivity and ordering of the roots there, $p_{1} < 0$ at $0.27 > \tfrac{\pi}{12}$, and strict increase of $r$ on ten sample points of $(0, \tfrac{\pi}{12})$; with the monotonicity above, $\delta_{*}$ lies in the stated window. (Computed in floating point, for orientation only: $\delta_{*} \approx 0.2222220471$; the Koide prediction of the tau mass at $\delta_{*}$ is $1776.969$ MeV, within the stored $1776.93 \pm 0.09$.)

\emph{(ii)} $\tfrac29 < \tfrac{\pi}{12}$ (the check uses the Machin lower bound on $\pi$), and the certified interval for $r(\tfrac29)$ lies in $(206.7703, 206.7704)$. Its lower end exceeds $r_{\max}$ by more than $400\sigma_{\mu}/m_{e}$, and $\sigma_{\mu}/m_{e}$ bounds the standard deviation of the measured ratio from above (the electron's relative uncertainty is $3 \times 10^{-10}$ against the muon's $2.2 \times 10^{-8}$). By monotonicity $\delta_{*} < \tfrac29$, and the check certifies $\tfrac29 - 0.22222204830 > 1.7 \times 10^{-7}$. (Computed in floating point, not part of the proof: with $\delta = \tfrac29$ held fixed and $\lambda$ solved from $r_{\mathrm{obs}}$, $\lambda^{2} = 1.99999915$, $Q - \tfrac23 = -1.4 \times 10^{-7}$ and the predicted tau mass is $1776.967$ MeV, inside the tau window. So $\tfrac29$ alone is not excluded by the data; the joint exactness of $Q = \tfrac23$ and $\delta = \tfrac29$ is. The value $\tfrac29$ for this phase is Brannen's observation, 2006, cited from memory.)

\emph{(iii)} The check finds, by direct search over denominators with integer arithmetic, the least $d$ for which some $k/d$ lies in $[0.22222204589, 0.22222204830]$: $d = 630131$, with $k/d = \tfrac{140029}{630131}$.

\emph{(iv)} $\pi$ is bracketed by Machin's formula $\pi = 16\arctan\tfrac15 - 4\arctan\tfrac{1}{239}$, each arctangent summed to $40$ terms with the first omitted term as error (alternating, decreasing), giving an interval of width below $10^{-50}$ inside $(3.14159265358979, 3.14159265358980)$. The least denominator of a rational in $[\delta_{\mathrm{lo}}/\pi_{\mathrm{hi}}, \delta_{\mathrm{hi}}/\pi_{\mathrm{lo}}]$ is $25249$ (first fraction $\tfrac{1786}{25249}$). If some orbit point $\pm\delta_{*} + \tfrac{j\pi}{6}$ equalled $\tfrac{a}{b}\pi$, then $\delta_{*}/\pi = \pm(\tfrac{a}{b} - \tfrac{j}{6})$ would be a rational of denominator dividing $6b$, so $6b \ge 25249$, $b \ge 4209$. Hence no point of the orbit, and so no channel phase $\chi$ on either slice under any relabelling, is a rational multiple of $\pi$ of denominator at most $4208$.

\emph{Reading of (iv) (argued).} A phase fixed by invariance under a symmetry of finite order $n$ is fixed by a root-of-unity condition and lies in $\tfrac{2\pi}{n}\mathbb{Z}$, possibly after the shifts of the orbit. The stabilisers met in face space are of order at most $6$: the $\mathbb{Z}_{2}$ left by the real structure (Proposition \claimref{synthesis.real-structure-breaks-source-stabiliser-to-z2}(ii)), the $S_{3}$, the quarter turn $\exp(\tfrac{I_4\pi}{2}P_{d})$ (numbers read in $\mathrm{span}\{1, I_4\}$) and the $\mathbb{Z}_{3}$ of phases with unit product that fixes the cubic family word (ibid.\ (iv)). None of them, nor any finite symmetry of order up to $4208$, can fix the phase that the measured split requires at exact Koide. This is consistent with, and independent of, the stored results that the real-structure completions do not see the channel phase at all and that the cubic word forces nothing.

\emph{(v)} With $\epsilon = 10^{-6}$ the certified difference quotient $(r(\delta(1 + \epsilon)) - r(\delta))/(r(\delta)\epsilon)$ lies in $(12.4, 12.5)$ at both ends of the window. (Floating point: $d\ln r/d\ln\delta = 12.48$ at $\delta_{*}$.) The relabelling signature of Proposition \claimref{families.classification-displaced-gaussian-class-solved-level} is unit sensitivity of the split to a free parameter; the channel phase does not carry it, so a rule fixing the phase would not be flagged by that condition, but it must fix the phase to a relative precision about $12.5$ times finer than the precision wanted for the split ($2.4 \times 10^{-9}$ absolute for the measured precision of $m_{\mu}/m_{e}$).

\emph{Scope.} The result assumes the projection reading of Proposition \claimref{synthesis.coherent-face-vector-real-slices} (the projected coherent vector read as the $\sqrt{m}$ vector) and exact Koide; it says nothing about states whose mass form is not of that form, and nothing about a rule that produces the Koide relation only approximately. It is transport-absent: data, the stored parametrisation and elementary number theory.
